\documentclass{article}
\usepackage{booktabs}
\usepackage{graphicx}
\usepackage{float}
\usepackage{geometry}
\geometry{a4paper, margin=1in}

\title{Analysis of Quantum Advantage Search Results}
\author{Automated Analysis Report}
\date{\today}

\begin{document}

\maketitle

\section{Abstract}
This report summarizes the findings from the quantum advantage search experiments conducted on January 26, 2026. We compare the performance of distinct quantum-classical hybrid architectures against classical baselines (ResNet50 and ViT). The results indicate that while classical baselines currently maintain superiority in both accuracy and training efficiency, specific quantum configurations—notably Amplitude Encoding with Variational Quantum Circuits (VQC)—demonstrate competitive performance potential.

\section{Methodology}
The experiments evaluated various combinations of:
\begin{itemize}
    \item \textbf{Backbones}: ResNet50, Vision Transformer (ViT).
    \item \textbf{Quantum Layers}: VQC (Amplitude, Angle, IQP encodings), QLSTM, QAOA.
    \item \textbf{Qubit Counts}: 0 (Classical), 4, 8.
    \item \textbf{Depths}: 2, 3.
\end{itemize}
Training metrics including best accuracy, final accuracy, final loss, and average epoch time were recorded.

\section{Results}
Table \ref{tab:results} presents the comparative metrics for the evaluated configurations.

\begin{table}[H]
    \centering
    \caption{Performance comparison of classical and quantum-hybrid architectures.}
    \label{tab:results}
    \begin{tabular}{lcccc}
        \toprule
        \textbf{Configuration} & \textbf{Best Acc} & \textbf{Final Acc} & \textbf{Final Loss} & \textbf{Avg Time (s)} \\
        \midrule
        Baseline ResNet50 (Classical) & 0.9611 & 0.9611 & 0.0679 & 25.12 \\
        Hybrid ResNet (Amp VQC, Q8, D2) & 0.9420 & 0.8883 & 0.6508 & 55.69 \\
        Hybrid ResNet (Amp VQC, Q8, D3) & 0.9290 & 0.9290 & 0.6916 & 67.17 \\
        Baseline ViT (Classical) & 0.9086 & 0.9086 & 0.2430 & 254.69 \\
        Hybrid ResNet (Angle VQC, Q8, D2) & 0.4975 & 0.2475 & 2.0524 & 61.01 \\
        Hybrid ViT (QLSTM, Q4, D2) & 0.2284 & 0.2235 & 1.9060 & 415.41 \\
        Hybrid ResNet (IQP VQC, Q8, D2) & 0.2191 & 0.0901 & 2.2903 & 84.03 \\
        Hybrid ViT (IQP QAOA, Q4, D2) & 0.1346 & 0.0599 & 2.2987 & 490.14 \\
        \bottomrule
    \end{tabular}
\end{table}

\section{Discussion}
\subsection{Classical vs. Quantum Performance}
The classical ResNet50 baseline achieved the highest accuracy (96.11\%) and the lowest computational cost (25.12s/epoch). This establishes a strong benchmark that current hybrid approaches strive to match. The classical ViT, while reaching decent accuracy (90.86\%), incurred a significant computational penalty (254.69s/epoch), rendering it less efficient than the ResNet architecture in this study.

\subsection{Efficacy of Amplitude Encoding}
Among the quantum variants, the \textit{Hybrid ResNet with Amplitude VQC} (8 qubits) showed remarkable promise, achieving a best accuracy of 94.20\%. This result is comparable to the classical baseline, suggesting that amplitude encoding effectively preserves relevant feature information within the quantum circuit. However, deep circuits (Depth 3) did not yield improvements over shallower ones (Depth 2), potentially indicating trainability challenges such as barren plateaus or overfitting.

\subsection{Encoding Strategy Impact}
The choice of encoding strategy proved critical. While Amplitude Encoding succeeded, Angle Encoding and IQP-based methods resulted in significantly degraded performance (Accuracy $< 50\%$). This disparity suggests that for the given dataset and task, the dense information packing of amplitude encoding is superior to the feature mapping capabilities of angle or IQP encodings.

\subsection{Architecture Latency}
Quantum hybrids consistently showed higher training latencies than their classical counterparts. For instance, the Hybrid ResNet (Amp VQC) took approximately $2.2\times$ longer per epoch than the classical ResNet. The ViT-based hybrids were exceptionally computationally intensive, with epoch times exceeding 400 seconds.

\section{Conclusion}
The study highlights that while quantum advantage in terms of raw accuracy or speed is not yet realized in this specific setup, hybrid architectures utilizing Amplitude Encoding demonstrate viability, approaching classical performance levels. Future work should focus on optimizing quantum circuit depth and exploring more efficient training protocols to bridge the remaining performance gap.

\end{document}
